\documentclass{article}

\usepackage{amsmath,amssymb}
\usepackage{xurl}
\usepackage[colorlinks=true,linkcolor=blue,citecolor=blue,urlcolor=blue]{hyperref}
\setlength{\emergencystretch}{2em}

% Shared commands for notes in docs/paper-gaps/.

\DeclareUnicodeCharacter{2080}{\ifmmode{}_{0}\else\textsubscript{0}\fi}
\DeclareUnicodeCharacter{2081}{\ifmmode{}_{1}\else\textsubscript{1}\fi}
\DeclareUnicodeCharacter{2082}{\ifmmode{}_{2}\else\textsubscript{2}\fi}
\DeclareUnicodeCharacter{2099}{\ifmmode{}_{n}\else\textsubscript{n}\fi}

\newcommand{\Lean}{\textsc{Lean}}
\ExplSyntaxOn
\NewDocumentCommand{\leanid}{m}{
  \ifmmode
    \textsf{\detokenize{#1}}
  \else
    \group_begin:
    \urlstyle{sf}
    \tl_set:Nn \l_tmpa_tl {#1}
    \tl_replace_all:Nnn \l_tmpa_tl {\_} {_}
    \exp_args:NV \path \l_tmpa_tl
    \group_end:
  \fi
}
\ExplSyntaxOff
\providecommand{\tr}{\operatorname{tr}}
\newcommand{\tnleanIssueBase}{https://github.com/LionSR/TNLean/issues/}
\newcommand{\tnleanPrBase}{https://github.com/LionSR/TNLean/pull/}
\newcommand{\tnleanDocBase}{https://sirui-lu.com/TNLean/docs/}
\newcommand{\ghissue}[1]{\href{\tnleanIssueBase#1}{GitHub issue \##1}}
\newcommand{\ghpr}[1]{\href{\tnleanPrBase#1}{PR \##1}}
\newcommand{\doclink}[1]{\href{\tnleanDocBase#1.html}{\path{#1}}}

% Dirac notation and the identity operator, as in blueprint/src/macros/common.tex.
% \providecommand keeps note-local definitions authoritative.
\usepackage{dsfont}
\providecommand{\ket}[1]{|#1\rangle}
\providecommand{\bra}[1]{\langle #1|}
\providecommand{\braket}[2]{\langle #1 | #2 \rangle}
\providecommand{\ketbra}[2]{|#1\rangle\!\langle #2|}
\providecommand{\Id}{\mathds{1}}
\providecommand{\one}{\mathds{1}}
\providecommand{\C}{\mathbb{C}}
\providecommand{\MN}[1]{M_{#1}(\C)}
\providecommand{\MD}{M_D(\C)}

% Verdict marker: \gapnote{<kind>}{<status>}. Kind is the policy
% classification (clarification, local-correction, scope-restriction,
% unfaithful, false-source, open-gap); status is open, wip (elimination
% underway), resolved, or historical. Machine-read by the site index;
% prints a verdict line.
\newcommand{\gapnote}[2]{\par\noindent\textbf{Verdict:} #1 (#2).\par}


% Fallbacks keep this author document compilable when it is built outside its
% source directory and a different command.tex is found first.
\providecommand{\leanid}[1]{\textsf{\detokenize{#1}}}
\providecommand{\doclink}[1]{%
    \href{https://sirui-lu.com/TNLean/docs/#1.html}{\path{#1}}}
\providecommand{\gapnote}[2]{%
    \par\noindent\textbf{Verdict:} #1 (#2).\par}
\providecommand{\norm}[1]{\lVert #1\rVert}
\providecommand{\ket}[1]{|#1\rangle}
\providecommand{\bra}[1]{\langle #1|}

\title{The Frustration-Free Sweep by Range Projections}
\author{}
\date{2026-11-20}

\begin{document}
\maketitle
\gapnote{scope-restriction}{resolved}

\section{The source sweep}

The finite-range frustration-free specialization of the two-dimensional
area law of September~24, 2026~\cite{OpenAI2026AreaLaw}, recorded in the
reading note of November~18, 2026
(\path{docs/audits/2026-11-18_finite_range_frustration_free_area_law.tex},
lines~122--131), replaces the Poisson-clock constraints of the source's
Proposition~10.2 by an exact sweep of local filters.

Fix the model hypotheses of that note: a finite lattice domain
$\Lambda\subset\mathbb Z^2$, local positive terms $0\le h_X\le J\cdot\mathbf 1$
with $h_X\Omega=0$, a unit ground vector $\Omega$ with projection
$P=\ket\Omega\bra\Omega$, and the global gap
$\Delta(\mathbf 1-P)\le\sum_X h_X$. Write $k_i=h_i/J$, colour the intersection
graph of the supports in $\chi\le d+1$ colours with degree bound $d$, and set
\begin{align}
    G_i    & =(\mathbf 1-k_i)^{1/2}, &
    K_i    & =k_i^{1/2},             &
    D      & =G_N\cdots G_1 .
    \label{eq:sweep_source_filters}
\end{align}
The note's detectability argument gives, for $v\perp\Omega$,
\begin{align}
    \norm{D^t(\mathbf 1-P)} & \le\theta^t, &
    \theta                  & =\bigl(1+g/(d+1)^2\bigr)^{-1/2}<1, &
    g                       & =\Delta/J.
    \label{eq:sweep_source_bound}
\end{align}

\section{The formal deviation}

The Lean development
(\path{TNLean/PEPS/AreaLaw/FrustrationFree/SweepContraction.lean} and
\path{TNLean/PEPS/AreaLaw/FrustrationFree/RangeProjection.lean}) does not use
the square-root filters~\eqref{eq:sweep_source_filters}. It uses instead the
orthogonal projection $\Pi_i$ onto the range of $h_i$---equivalently onto
$(\ker h_i)^\perp$, the support projection of the positive semidefinite
$h_i$---and the complementary factors $\mathbf 1-\Pi_i$, with
$D=\prod_i(\mathbf 1-\Pi_i)$ in the same colour order.

This changes the sweep operator but no hypothesis of the model:
\begin{itemize}
    \item $\Pi_i$ has the same local support as $h_i$, because a support
          projection is a continuous-functional-calculus value of $h_i$.
    \item $h_i\Omega=0$ gives $\Pi_i\Omega=0$, so each factor $\mathbf 1-\Pi_i$
          fixes $\Omega$ and commutes with $P$, exactly as $G_i$ does.
    \item Since $0\le k_i\le\mathbf 1$ and $k_i$ vanishes on $\ker\Pi_i$, one
          has $k_i\le\Pi_i$; hence
          $\sum_i\Pi_i\ge\sum_i k_i=J^{-1}\sum_i h_i\ge g(\mathbf 1-P)$, the same
          gap input with $g=\Delta/J$.
\end{itemize}
The norm-loss argument of QICLean's detectability lemma applies to the
positive contractions $\mathbf 1-\Pi_i$ unchanged, because a projection's own
contribution vanishes once its factor has acted. It therefore reproduces the
bound~\eqref{eq:sweep_source_bound} with the note's constant $\theta$ (and in
fact admits the sharper denominator $d^2$ in place of $(d+1)^2$, so the formal
constant is at least as good as the source's). No hypothesis of the model is
added, and the downstream light-cone, amplification, and two-family-assembly
steps use~\eqref{eq:sweep_source_bound} verbatim.

\section{Scope}

The deviation is confined to the choice of sweep filters. The formal
\eqref{eq:sweep_source_bound} is the note's sweep contraction, established by a
mathematically equivalent route that is correct as stated and introduces no
side hypothesis. The restriction to projection filters is therefore resolved:
it is an alternative derivation of the same inequality, not a weaker theorem.
The formalization plan is \path{docs/formalization/frustration-free-area-law-plan.md}.

\bibliographystyle{alpha}
\bibliography{references}

\end{document}
